\chapter{Appendix}

The appendix holds the reference material that supports the main survey. It has three parts.

\begin{itemize}
  \item A \textbf{comparative table of terms}, one plain-English concept per row and the nearest term in each
  of the six core traditions, for quick cross-reading.
  \item The \textbf{per-tradition term mappings}, which take each tradition in turn and give its key native
  terms with their plain-English meanings, with native script where it adds value and a standard romanization
  first.
  \item The \textbf{creation stories in native terms}, each tradition's creation account retold once more,
  this time using its own vocabulary and names, as a direct parallel to the plain-English creation model
  given in the body.
\end{itemize}

The body of the survey keeps to plain English so the structures show through. The appendix restores the
native vocabulary, so a reader can move between the bare model and its traditional dress.
